\chapter{Requirements Specification}
\label{ch:req}

\section{Functional Requirements}

\begin{enumerate}[label=\textbf{FR\arabic*.}, widest=10, leftmargin=*]
  \item The system shall connect to a target AI model through its text API, send
  prompts, and read responses.
  \item The red agent shall select attack prompts from a prepared set and/or
  generate new ones, covering known attack categories.
  \item The blue agent shall detect and block unsafe responses using guardrails
  and output filters, and shall label each response as safe or unsafe.
  \item The system shall collect evidence for every attack: the prompt, the
  model's response, whether the model refused or complied, and any guardrail
  signal.
  \item The dual-reasoning engine shall produce an evidence-first conclusion and a
  hypothesis-first conclusion, and shall accept a result only when both agree.
  Otherwise, it shall mark the result as low-confidence and flag it for a human.
  \item For a missed attack, the reflection step shall find the root cause and
  propose a concrete, testable fix.
  \item The validation step shall re-test each fix to confirm that it blocks the
  attack.
  \item The system shall store evidence, reasoning path, confidence, outcome, and
  lessons in an auditable memory.
  \item The learning loop shall apply verified lessons to the next cycle.
  \item Human approval shall be required before deploying fixes or other
  high-impact changes.
\end{enumerate}

\section{Non-Functional Requirements}

\begin{itemize}
  \item \textbf{Performance:} the system should complete one
  attack--reason--validate cycle in a reasonable time.
  \item \textbf{Accuracy:} dual-path reasoning should reduce wrong conclusions
  compared with a single reasoning path.
  \item \textbf{Reliability:} the system should run repeated cycles without
  losing memory or state.
  \item \textbf{Usability:} the system should produce clear logs and summaries of
  what was attacked, what was caught, and what was fixed.
  \item \textbf{Scalability and Extensibility:} the system should work with any
  target that exposes a text API, and it should be easy to add new attack
  categories.
  \item \textbf{Security and Privacy:} the system should operate only on an
  authorized target, use no real user data, and keep all testing inside a
  controlled environment.
\end{itemize}

\section{Hardware Requirements}

\begin{itemize}
  \item Multi-core CPU and 16\,GB or more of RAM.
  \item Optional GPU. This helps when running a local open-source model, but it
  is not required when using a cloud-hosted model through an API.
  \item Sufficient storage (an SSD is recommended).
  \item A reliable internet connection (required when using a cloud-hosted
  model).
\end{itemize}

\section{Software Requirements}

\begin{itemize}
  \item Operating system: Windows or Linux.
  \item Python 3.10 or higher.
  \item Standard AI/agent and HTTP libraries.
  \item Open-source red-teaming and attack-generation tools.
  \item An open-source guardrail / safety-classifier toolkit (for the blue side).
  \item Access to an open-source LLM, run locally or through a cloud-hosted
  inference API.
\end{itemize}

\section{Constraints and Assumptions}

\subsection{Constraints}
\begin{enumerate}
  \item Phase~1 targets AI systems (LLM, chatbot, or agent) that are reached
  through a text API.
  \item Attacks are performed only on a target that the team controls and is
  authorized to test --- never on live third-party services.
  \item The project uses pre-trained open-source models. It does not train a
  foundation model from scratch.
  \item Evaluation uses a limited set of attack categories and prompts because of
  time and compute limits.
\end{enumerate}

\subsection{Assumptions}
\begin{enumerate}
  \item The target model exposes a text input/output API.
  \item The attack categories drawn from public sources reasonably cover current
  threats.
  \item A human reviewer is available to approve high-impact changes.
\end{enumerate}
